\subsection{Confirmatory study}\label{sec:confirmatory}

The four hypotheses of the prespecified protocol were evaluated on 53 new instances (Table~\ref{tab:confirmatory}). Every decision uses the unrounded interval limits. The protocol's validation passed before any quantity was computed.

\begin{table}[htbp]
\centering
\footnotesize
\caption{Confirmatory results on 53 fresh instances. Differences are in percentage points of validated service-event probability; intervals are paired bootstrap intervals over instances. Conservative speed returned a plan on 50 of 53 instances, and H1 and H4 use those instances as the protocol defines.}
\label{tab:confirmatory}
\begin{tabular}{@{}>{\raggedright\arraybackslash}p{1.9cm}>{\raggedright\arraybackslash}p{4.2cm}>{\raggedright\arraybackslash}p{3.5cm}cc@{}}
\toprule
Hypothesis & Quantity & Estimate [interval] & Level & Decision\\
\midrule
H1 (primary) & Original procedure minus conservative speed, 50 instances & 15.7 [7.5, 23.6] & 97.5\% & confirmed\\
H2 (primary) & Capped minus original procedure, 53 instances & 11.7 [6.0, 17.9] & 97.5\% & confirmed\\
H3 & Admitted by capped / original procedure & 46 / 32 (capped only 14, original only 0) & exact $p=0.00012$ & confirmed\\
H4 & Advantage at $\sigma=0.30$ minus at $\sigma=0.10$ ($\rho=0$), 50 instances & 28.3 [19.6, 36.4] & 95\% & confirmed\\
\bottomrule
\end{tabular}
\end{table}

All four hypotheses were confirmed. The original procedure exceeded conservative speed by 15.7 points, capping exceeded the original procedure by 11.7 points, and capping admitted 46 instances against 32. Capping lost no instance that the original procedure admitted (14 gained, 0 lost). The advantage over conservative speed was 28.3 points larger at $\sigma=0.30$ than at $\sigma=0.10$ for independent disturbances. Of the plans admitted on the screening bank, 31 of 32 (original) and 45 of 46 (capped) retained the 0.95 lower bound on the fresh bank. Pooled over all instances, service probability was 79.4\% for the original procedure, 82.4\% for screened conservative speed and 91.0\% for the capped procedure, so the 95\% target is still not met on average. Screened conservative speed was not a prespecified comparison, and its pooled value is descriptive only. As in the development study, the shortfall comes from instances without an admitted plan. Plans admitted by screening averaged 97.4\% for the original procedure and 97.5\% for the capped procedure, while the fallback plans averaged 51.9\% (21 instances) and 48.7\% (7 instances).

H2 and H3 are not independent. When neither procedure admits a plan, both return the same fallback, and 11.6 of the 11.7 points of H2 come from the 14 instances that only the capped procedure admits. The two results are one finding measured twice.

\paragraph{Post hoc fallback.} The protocol fixes the fallback: without an admitted plan, both procedures return the uncontracted slack-aware plan. A variant that returns the conservative-speed plan instead, where that plan exists, would have raised pooled service on the same fresh bank from 79.4\% to 85.0\% for the original procedure (18 instances changed) and from 91.0\% to 93.0\% for the capped procedure (6 instances changed). The fallback rule uses no validation outcome, but this variant was chosen after the results were seen, so it is exploratory.

These instances come from the same generator and the same assumed disturbance model as the development study. The result establishes repeatability within that setting. It does not address real road networks, calibrated parameters or other instance families.
